\documentclass{book}

\usepackage{lipsum}
\usepackage{amsmath}
\usepackage{hyperref}
\usepackage[capitalize, noabbrev]{cleveref}

\begin{document}
	% First Chapter
	\chapter{Introduction}
	\label{chap:intro}
	\lipsum[1] \\
	
	Let's refer to \cref*{chap:ins} here !!
	
	\section{Basic Concepts}
	\label{sec:basicons}
	This is a range reference for 
	\Crefrange{eq:1steq}{eq:3rdeq} \\
	
	This is a bulk reference of all labeled objects
	in the last chapter
\cref{eq:2ndeq,eq:1steq,eq:3rdeq,chap:conc,sec:fineqs}

\cref{eq:1steq,eq:2ndeq,,eq:3rdeq,eq:frtheq,eq:fiftheq}
	
	% Chapter Two
	\chapter{Insights}
	\label{chap:ins}
	\lipsum[1] \\
	
	This is a reference to \cref{sec:basicons} \\
	
	\Cref{eq:1steq} has its own reference capitalized
	and not abbreviated unlike this one
	\cref{eq:1steq} in the middle of a sentence.
	
	% Chapter Three
	\chapter{Conclusion}
	\label{chap:conc}
	\lipsum[1] 
	
	\section{Final Equations}
	\label{sec:fineqs}
	
	\begin{equation}
		F(x) = 2x + y
		\label{eq:1steq}
	\end{equation}
	\begin{equation}
		F(x) = 2x + y
		\label{eq:2ndeq}
	\end{equation}
	\begin{equation}
		F(x) = 2x + y
		\label{eq:3rdeq}
	\end{equation}
	\begin{equation}
		F(x) = 2x + y
		\label{eq:frtheq}
	\end{equation}
	\begin{equation}
		F(x) = 2x + y
		\label{eq:fiftheq}
	\end{equation}
\end{document}